% State of the Art -- chapter opening. Paste into your own thesis, or \input it.
% Section files in this folder are meant to be included in numeric order.
% Follows the Clinical Background chapter: CAD, CCTA, WSS, FFR and the artery names are already
% defined there and are not redefined here.

\chapter{State of the Art}
\label{ch:sota}

Objective 4 asks for two things: the state of the art in segmenting the coronary arteries, and in
describing their topology.
This chapter answers them in that order, then narrows to where the second one stops.

Section~\ref{sec:segmentation} covers the segmentation half: the methods ImageCAS-X benchmarks
against each other and against inter-observer agreement, and why this thesis builds its framework
on CAS-Net.
Section~\ref{sec:scale-gap} opens the topological-description half with the evidence matrix the
rest of the chapter is read against: where the literature linking coronary shape to disease actually
sits, across three scales.
Sections~\ref{sec:getting-tree} and~\ref{sec:describing-tree} cover what is already possible: a tree
can be built from a segmentation reliably enough to measure, and a population of trees has already
been described, twice, in cohorts of a few hundred.
Sections~\ref{sec:geometry} and~\ref{sec:shape-enough} cover what tree-scale geometry is already
known to predict, feature by feature, and why that mechanism does not require simulating blood flow
to use.
Section~\ref{sec:nearest-competitors} places this thesis against the work that sits closest to it
without doing what it does.
Section~\ref{sec:the-gap} states what none of the above has done: measured tree-scale geometry across
a population of thousands, described its distribution, and asked which individuals fall outside it.
That is the gap objectives 6 through 9 occupy.
